\scanpage{108}
% archive sheet number: 107
\uline{Ex}\quad $N=3$

\textcircled{1}\quad \fbox{$A^2(X)$ fini} \quad en particulier $J^{\unsure{2|1}}(X)=0$

\textcircled{2}\quad \fbox{$J^1(X)\longrightarrow J^1(Y)$ \quad injectif exc. $p$}\deleted{\unsure{injectif}}

\textcircled{3}\quad \deleted{$\unsure{A^2(X)|\mathcal{A}^2(X)}$} contenu dans 1º)

groupes essentiels \quad $J^1(X),N^1(X)$ \qquad définissant $A^1(X)$

\unsure{hom|homs} essentiels
\[
\begin{aligned}
 &J^1(X)\longrightarrow J^1(Y)\\
 &N^1(X)'\longrightarrow J^1(Y)
\end{aligned}\ \bigg|\ A^1(X)'\longrightarrow J^1(Y)
\]
% slanted
\margin{sous-gpe d'indice fini de $N^1(X)$}

\uline{Pb} définir un complémentaire naturel de $J^1(X)$ dans $J^1(Y)$ (correspondant aux ``cycles \unsure{évanescents|évanescent}'')

Donner un sens à la \unsure{formule}
\[
 J^1(X)\overset{\sim}{\longrightarrow}J^1(Y)^{\pi}
\]
Interpréter le th. de \inserted{l'index de} Hodge pour les surfaces projectives en termes de la négativité d'une certaine forme définie sur $N^1(X)$\ldots?

\uline{Ex}\quad $N=4$

\textcircled{1}
\[
 \left\{\begin{aligned}
 &N^2(X)\ \text{fini}\\
 &A^3(X)\ \text{fini, en particulier }J^3(X)=0
 \end{aligned}\right.
\]
\textcircled{2}
\[
\begin{aligned}
 J^1(X)&\longrightarrow J^1(Y)&&\text{bij exc. $p$}\\
 \deleted{\unsure{J}}N^1(X)&\longrightarrow N^1(Y)&&\text{inj. exc $p$}\\
 &\Rightarrow\quad A^1(X)\longrightarrow A^1(Y)&&\text{\uline{bij exc.} $p$}
\end{aligned}
\]
\deleted{\textcircled{3}}
